\chapter{分子的形状决定它会做什么}

\begin{kaichang}
剥开一只橙子，指尖沾上一点果皮油，凑近闻一闻：清新的橙子香，几乎人人都熟悉。再剥一只柠檬，同样沾一点果皮油——气味立刻换了频道，酸冽、明亮，绝不会和橙子弄混。

现在告诉你一件怪事。橙子和柠檬果皮油里最主要的香味分子，由完全相同数量的原子组成：10 个碳、16 个氢。不只如此，原子之间的连接顺序也一模一样。写出化学式，都是 \chem{C_{10}H_{16}}；画出“哪个原子连着哪个原子”的连接图，两张图完全重合。可是它们的气味就是不同，一个让你想到橙子，一个让你想到柠檬。

这两个分子的差别到底藏在哪里？先猜一猜，写在纸上。本章结束时揭晓答案。线索只有一句：差别不在“零件清单”上，而在三维空间里。
\end{kaichang}

\section{先闻为快：形状是能被“摸”出来的}

我们看不见单个分子，但生活早就反复提醒你：分子的\textbf{形状}是有后果的。

水和干冰是最好的一对。水和二氧化碳都含氧，组成都简单得不能再简单，可在室温下一个是液体，一个是气体——干冰甚至不屑于熔化，直接从固体跳成气体，在舞台烟雾里翻滚。还有：糖丢进水里会溶解，沙子丢进去纹丝不动；酒精擦在手背上挥发时凉凉的；一粒药能治病，换成分子式几乎一样的另一粒却可能毫无效果，甚至有害。

还有更直接的证据来自你的鼻子和皮肤。衣柜里的樟脑丸，分子能穿过塑料袋的缝隙飘出来，让整柜衣服染上气味——但同样放在柜子里的铁钉，你什么也闻不到。樟脑分子恰好是能挥发、又恰好能插进你鼻腔里某个受体口袋的形状；铁原子组成的晶体则牢牢抱在一起，根本不打算起飞。夏天用的防晒霜也是一场精心设计的形状工程：那类分子的骨架形状决定了它们恰好能吸收紫外线的能量，把伤害皮肤的辐射转成无害的热。形状合适的分子被吸收、被闻到、被利用；形状不合适的，从你的感官和身体旁边悄悄滑过，仿佛不存在。

这些现象的差别，用“由什么原子组成”解释不通，用“原子怎么连接”也常常解释不通。解释的钥匙藏在一个我们还没认真看过的维度上：分子不是平面图，它是三维空间里的一团东西。原子各占一块位置，电子云各占一块空间，整体呈现出一个确定的\textbf{几何形状}。而这一章要讲的规则出人意料地朴素：\textbf{负电荷互相排斥，所以中心原子周围的电子对会彼此推开，直到离得尽可能远}。分子的大部分形状，就是这条朴素排斥规则的几何后果。

\section{电子对为什么要把彼此推远}

把镜头推进到第 19 章认识的共价分子。中心原子周围有两类电子云：一类是被两个原子共享的\textbf{成键电子对}（双键、三键先粗略地各算“一团”电子云），另一类是只属中心原子自己、不参与成键的\textbf{孤对电子}。不管是哪一类，它们都是带负电的电子云，而负电与负电之间只有一件事：排斥。

排斥会怎样安排它们的位置？设想你站在房间中央，周围要放几个都讨厌彼此的吵闹气球，每个气球都想离别的气球越远越好：

\begin{itemize}[itemsep=2pt]
\item 2 个气球：只能一头一个，排成\textbf{直线}，夹角 $180^{\circ}$。
\item 3 个气球：在同一平面上互成 $120^{\circ}$，排成\textbf{平面三角形}。
\item 4 个气球：平面上最多挤下 3 个，第 4 个会把大家撑开成立体结构——正\textbf{四面体}的四个顶点，两两夹角约 $109.5^{\circ}$。
\end{itemize}

化学家给这套思路起了个名字：价层电子对互斥模型（VSEPR，念作“维斯珀”）。它本质上是第 17 章那条总原则的推论：体系会落在能量较低的状态，而电子对离得远，排斥的能量就低。分子没有眼睛、没有意愿，是静电排斥把形状“推”出来的。

\begin{shiyan}
\textbf{气球分子模型（在家可做）}

材料：同型号长气球 4～6 个、细绳或橡皮筋若干。

步骤：把气球吹成差不多大小的长条形（不要吹太满），每个气球打一个结。取 2 个气球，在中间段用细绳紧紧扎在一起，松手观察它们自动摆成什么姿势。再取 3 个、4 个，分别在中点扎成一束，观察自动形成的排列。

观察点：2 个气球是否自然排成直线？3 个是否自动摊成平面“奔驰车标”？4 个是否能躺平——还是会翘成一个立体的“三脚架加一竖”？最后，用一只手用力捏住 4 气球束中的一个气球（它代表一对“孤对电子”），观察其余三个气球的夹角是不是被挤小了。

安全提示：气球可能突然爆裂，扎结时远离眼睛；气球碎屑及时收走，避免幼童或宠物误食。

这个模型的局限也要说清楚：气球只模拟“电子云互相排斥、各占一块空间”，真实分子里没有橡皮膜，也没有颜色。
\end{shiyan}

\begin{qiaoliang}
为什么四面体的夹角偏偏是 $109.5^{\circ}$，而不是 $90^{\circ}$ 或 $120^{\circ}$？因为 4 个彼此排斥的点在球面上离得最远的摆法只有一个：各占正四面体的一个顶点。用简单的几何可以算出，从正四面体中心看任意两个顶点，夹角都是 $\arccos(-1/3) \approx 109.5^{\circ}$。

再看两个被“挤小”的数字。氨 \chem{NH_3}：中心氮周围也是 4 对电子，但其中 1 对是孤对电子。孤对电子只受一个原子核吸引，电子云摊得更“胖”、排斥更强，把 3 个成键对往下压，H–N–H 键角从 $109.5^{\circ}$ 被压到约 $107^{\circ}$。水 \chem{H_2O} 有 2 对孤对电子，挤压加倍，H–O–H 键角只剩约 $104.5^{\circ}$。这三个数字——$109.5^{\circ}$、$107^{\circ}$、$104.5^{\circ}$——不用背，记住“孤对电子更霸道，多一对就再挤小一点”就够了。
\end{qiaoliang}

\section{四种基本形状和它们的“名片”}

把 VSEPR 用到第 19 章的路易斯结构上，只需三步：数中心原子周围的电子对总数；分清几对成键、几对孤对；让电子对离得最远，然后\textbf{只看原子的位置}给分子形状命名。

拿水完整走一遍这三步。氧在最外层有 6 个电子（第 10 章的“楼层”知识：氧在第二行第 16 族），拿出 2 个分别和两个氢共享，形成 2 对成键电子；剩下 4 个电子配成 2 对孤对。总数 4 对——所以电子对的几何骨架是四面体。再看原子落在哪儿：4 个顶点里 2 个被氢原子占据，2 个被“隐身”的孤对占据，于是我们看到的分子是折线形。同一份推演换个原子：氮最外层 5 个电子，成 3 对键、剩 1 对孤对，所以氨是三角锥；碳最外层 4 个电子，4 对全部成键，所以甲烷是完美的四面体。你发现了没有：只要会数最外层电子，分子的形状几乎可以\textbf{算}出来，而不是背出来。

常见的几张“名片”：

\begin{itemize}[itemsep=2pt]
\item \textbf{直线形}：二氧化碳 \chem{CO_2}。中心碳有 2 团电子云（两个双键），互成 $180^{\circ}$。
\item \textbf{平面三角形}：三氟化硼 \chem{BF_3}。中心硼周围 3 对成键电子，互成 $120^{\circ}$，整个分子是平的。
\item \textbf{四面体形}：甲烷 \chem{CH_4}。碳周围 4 对全是成键电子，4 个氢各占四面体一个顶点。
\item \textbf{三角锥形}：氨 \chem{NH_3}。电子对几何是四面体，但 1 个顶点被孤对电子占据，3 个氢原子和氮原子看起来是一座三角锥。
\item \textbf{折线形}：水 \chem{H_2O}。4 对电子中 2 对是孤对，两个氢原子在氧的同一侧，分子像字母 V。
\end{itemize}

注意最后一个细节：命名分子形状时，孤对电子“隐身”了——我们说水是折线形，不说它是四面体形，因为眼睛（和仪器）看到的是原子的位置。但决定这个形状的，恰恰是包括孤对在内的全部电子对。

顺便看一眼“名片册”的后面几页，知道存在即可：5 对电子离得最远的方式是\textbf{三角双锥}（3 个围成赤道的三角形，上下各 1 个，如 \chem{PCl_5}），6 对则是\textbf{八面体}（上下左右前后，如 \chem{SF_6}）。规律始终是同一条：先让电子对在中心原子周围彼此离得最远，再看原子落在哪儿。

形状到了纸面上还需要一套符号。化学家画立体分子有一套约定：普通直线表示键在纸面上，\textbf{实心的楔形键}表示这个原子朝你伸出纸面，\textbf{虚线键}表示它缩到纸面背后。于是甲烷不再画成平面十字——那是错误的暗示——而是一个碳连着两个纸面上的氢、一个楔形键指向前方的氢、一个虚线键指向后方的氢。这套符号你将在后面每一章反复见到。

\begin{tikzpicture}[scale=1.1]
\fill[gray!55] (0,0) circle (0.32);
\node at (0,-0.02) {\chem{O}};
\draw[thick] (0.22,0.22) -- (0.72,0.72);
\draw[thick] (-0.22,0.22) -- (-0.72,0.72);
\fill[white] (0.86,0.86) circle (0.2);
\draw (0.86,0.86) circle (0.2);
\node at (0.86,0.84) {\chem{H}};
\fill[white] (-0.86,0.86) circle (0.2);
\draw (-0.86,0.86) circle (0.2);
\node at (-0.86,0.84) {\chem{H}};
\fill[gray!30] (0.28,-0.3) circle (0.14);
\fill[gray!30] (-0.28,-0.3) circle (0.14);
\draw[dashed] (0.55,0.95) arc (40:140:0.75);
\node at (0,1.45) {$104.5^{\circ}$};
\node[align=center] at (0,-0.85) {两个灰点代表孤对电子（示意）};
\end{tikzpicture}

上图是水分子的折线形示意：这是人为的表示法，圆圈大小和灰度不代表真实外观，孤对电子也不是两颗小珠，只是两团占据空间的电子云。

\section{都是极性键，为什么水和二氧化碳不一样}

现在可以兑现本章标题的一半了：形状决定性质。第一个证据，就是水和二氧化碳这对“都含极性键、命运却相反”的分子。

回忆第 19 章：氧的电负性大，共享电子被拉向氧，所以 C=O 键和 O–H 键都是\textbf{极性键}——键的一端偏负、一端偏正。但“每根键有极性”和“整个分子有极性”是两回事，中间隔着形状。

二氧化碳是直线形：\chem{O=C=O}。两个极性键大小相同、方向正好相反，就像两个人在拔河，势均力敌——电子云整体的“重心”就在正中央，正负抵消。所以 \chem{CO_2} 分子\textbf{整体没有偶极}，分子与分子之间几乎不粘，室温下自然是气体，固体干冰稍一受热就直接升华。

水是折线形：两个 O–H 键的极性不但不能抵消，还合伙在氧的一端堆出负电、在氢的一端露出正电。每个水分子都是一个小小的偶极子，一头偏负、一头偏正。后果极其深远：水分子会彼此拉扯、会抱住很多溶质粒子、会把冰撑成镂空结构——这些精彩的连锁反应是第 22、23 章的主题，这里只记住因果链的起点：\textbf{极性键 + 不对称的形状 = 极性分子}。

再看一个容易被形状“骗”到的例子：四氯化碳 \chem{CCl_4}。C–Cl 键是极性键，但 4 个氯对称地占据四面体顶点，四股“拔河”的力恰好抵消，\chem{CCl_4} 整体没有偶极，所以它和水互不溶解，却能溶解油污。形状把键的极性“藏”了起来。

用同样的逻辑检查两个熟悉的分子。氨是三角锥形，孤对电子把分子撑出明显的一头偏负、一头偏正，所以氨是极性分子——这正是氨气极易溶于水、能制成氨水和化肥溶液的原因之一。甲烷 \chem{CH_4} 则是完美对称的四面体，四根 C–H 键的微弱极性全部抵消，所以天然气在水里几乎不溶。同样是“小分子”，溶解性格天差地别，判决书就是形状。

\section{形状是怎样被“数”出来的}

看不见的东西，怎么知道它长什么样？这是科学最迷人的部分。

\begin{zhengju}
1874 年，化学家面对一个难题：碳永远伸出 4 个键，可这 4 个键在空间怎么摆？谁也看不见。荷兰青年范特霍夫和法国青年勒贝尔（两人互不相识）用同一条推理破解了它——不是看，而是\textbf{数}。

他们的问题很具体：如果碳的 4 个键是平面的，比如指向正方形的四个角，那么二氯甲烷 \chem{CH_2Cl_2} 应该有两种版本——两个氯相邻（在正方形的邻角）和两个氯相对（在对角），就像同一副零件能拼出两样东西。可化学家们翻来覆去地制备、提纯、比较，全世界只找到\textbf{一种} \chem{CH_2Cl_2}。范特霍夫提出：只有当碳的 4 个键指向\textbf{四面体的四个顶点}时，任意两个位置才必然相邻，无论怎样交换都拼不出第二种——与事实吻合。

这个模型还顺手解开了一个 26 年前的悬案。1848 年，巴斯德用镊子在显微镜下把酒石酸盐的晶体分成两堆——两堆晶体的形状互为镜像，像左手套和右手套；溶于水后，一堆让偏振光向左转，一堆向右转。当时没人解释得了分子层面的原因。四面体碳给出预言：只要一个碳连着 4 个\textbf{不同}的基团，这个分子就可能存在一对镜像版本，化学性质几乎全同，唯有“旋光”和“与另一只手的匹配”不同。酒石酸正是这样的分子。

当然也有反对声。老牌化学权威科尔贝公开嘲讽范特霍夫是在“追逐幻影”，把想象的空间结构当成科学。裁决来自更硬的证据：20 世纪初，X 射线晶体学发展起来，科学家用 X 射线穿过晶体，从衍射图样反推出原子在空间中的真实位置——金刚石里的碳，四个键的夹角正是 $109.5^{\circ}$。一个靠“数异构体”猜出来的形状，被直接测量证实了。值得记住的方法：好的结构模型不靠看见，靠做出能被检验的预言。
\end{zhengju}

\section{这个模型什么时候会失灵}

VSEPR 是初中到大学都爱不释手的工具，但它是\textbf{经验模型}，有明确的适用范围。

它对主族元素的小分子（中心原子周围 2～6 对电子）几乎百发百中；遇到过渡金属的配合物就常常猜错——那些原子的 d 轨道（第 9 章提过的电子“户型”）参与成键，排斥规则不再简单。它告诉你电子对“想离远点”，却解释不了电子云真正的分布细节；更精细的图景要靠量子力学计算。遇到第 19 章说的共振结构（如臭氧 \chem{O_3}），也要先用路易斯结构数清楚电子对，VSEPR 才能开工。

还有一个比喻需要立刻拆穿：你也许听说过“杂化轨道”——碳的 1 个 s 轨道和 3 个 p 轨道“混合”成 4 个等价的 $\mathrm{sp^3}$ 轨道，指向四面体。请记住它的身份：杂化轨道是量子力学里为了方便描述而发明的\textbf{记账模型}，不是能被电子显微镜看见的实体。实验能测到的是整个分子的电子密度和键角，没有人“拍到”过一个杂化轨道。模型好用，不代表模型是实物——这是科学思维里极重要的一课。

最后，分子也不是雕像。化学键像弹簧，原子在平衡位置附近不停振动；单键还可以旋转，让分子的“姿势”不断改变。我们说“水的键角是 $104.5^{\circ}$”，说的是平均图景。对葡萄糖、蛋白质这样的大分子，整体形状由成千上万次折叠、扭转共同决定，单个原子的 VSEPR 只是拼图的一小块——第 48 章会专门讲蛋白质这台“折叠起来的分子机器”。

\begin{wuqu}
“课本插图上，氧原子是红球、氢原子是白球、键是小棍——分子就长这样。”

这是插图太逼真带来的误解。原子没有颜色（颜色来自分子对可见光的吸收，单个原子的“颜色”无从谈起）；化学键不是小棍，而是两团原子核之间概率较高的电子云；分子的边缘也不像气球皮那样界线分明，电子云是逐渐变淡的。球棍图、空间填充图、楔形虚线图，都是\textbf{为不同目的服务的表示法}：球棍图看连接和角度最清楚，空间填充图看分子“轮廓”最接近真实，但它们全是地图，不是领土。判断一张分子图好不好，标准不是“像不像”，而是“它想说明的问题有没有被说清”。
\end{wuqu}

\section{回到那颗橙子}

现在揭晓开场的悬念。橙子和柠檬果皮油里的主角是同一种分子——柠檬烯，\chem{C_{10}H_{16}}，连接方式完全相同。它们的差别，是\textbf{镜像}。

伸出你的两只手：手心朝上，大拇指都朝外。两只手的花纹、指头的连接顺序“完全相同”，但左手叠不到右手上——这就是\textbf{手性}。柠檬烯分子里有一个碳原子连着 4 个不同的基团，于是整个分子存在互为镜像、却不能重合的两个版本：一种主要散发橙子气味，另一种让你想到柠檬。薄荷与香菜籽（葛缕子）之间，也隔着同样一对镜像分子（香芹酮）。

为什么镜像会让鼻子分辨出来？因为你的嗅觉受体是蛋白质，是有着特定三维“口袋”的分子机器。气味分子飘进鼻腔，要插进口袋里才能触发神经信号。左手分子和右手分子插进同一只口袋，贴合的方式不同、触发信号的强度不同，大脑收到的“气味”就不同。因果链完整了：电子对排斥决定形状，形状决定它与受体怎样匹配，匹配决定你会闻出什么。\textbf{形状即功能}——这一原则将贯穿到全书最后一章。手性的故事远没讲完，第 42 章会专门细讲。

手性的后果不只关乎气味。20 世纪 50 年代末，药物沙利度胺作为镇静剂上市，分子也含有一个“手性中心”：一种镜像起镇静作用，另一种镜像会干扰胎儿发育，导致上万名婴儿出生缺陷，成为药学史上最惨痛的教训之一。后来的研究还发现，两种镜像在人体内会互相转化，所以即使只生产“好”的那一种也无法完全避险。这段历史的直接遗产是：今天的药物审批要求逐一研究药物每种镜像的作用。你药箱里的一半以上现代药物，都是只含单一镜像的“一只手”的分子。

形状即功能的原则，也被工业主动利用。香料工厂不再从成吨的橙皮里榨油，而是直接合成特定镜像的柠檬烯和香芹酮——只要形状对了，闻起来的“橙子”就是橙子。化肥工业的哈伯法把氮气和氢气变成氨，靠的也是催化剂表面原子的特定排列：氮分子躺在金属表面的“原子台阶”上，被形状合适的位点固定、削弱、拆开（第 30 章会细讲催化剂怎样另辟蹊径）。更日常的例子是洗洁精：它的分子一头亲水、一头亲油，正是这根“两头性格相反”的长条形状，让它能插进油滴和水之间，把油污从碗上撬下来冲走。分子世界里没有抽象的“功效”，一切功效都写在形状上。

镜像分子这么微妙，分子本身又有多小？做一个数量级估算：水分子宽约 $3\times 10^{-10}\,\mathrm{m}$（0.3 纳米）。一滴水约 0.05 克，除以水的摩尔质量 $18\,\mathrm{g/mol}$ 约为 $2.8\times 10^{-3}\,\mathrm{mol}$，乘上阿伏伽德罗常数 $\ud{6.02\times 10^{23}}{mol^{-1}}$，约含 $1.7\times 10^{21}$ 个水分子。把它们排成一列，长度约 $1.7\times 10^{21} \times 3\times 10^{-10}\,\mathrm{m} \approx 5\times 10^{11}\,\mathrm{m}$，即 5 亿千米——是地球到太阳距离的三倍多。一滴水里就藏着这么多“有形状的小东西”，而每一个的折线形状，都在为“水是液体”投票。

\section{形状之后，还剩下什么}

这一章从排斥规则出发，给分子拍了一张“三维证件照”：直线、平面三角、四面体、三角锥、折线。形状解释了水的极性、干冰的升华、橙子与柠檬的气味、药物的疗效与灾难。

但故事还缺一块。水分子是极性分子，\textbf{然后呢}？极性分子彼此靠近时会发生什么？偏负的一端遇到偏正的一端，会不会互相拉扯？正是这种拉扯，让水在室温下是液体而甲烷是气体，让壁虎能在天花板上行走，让露珠缩成小球。分子与分子之间这些“不粘成键、却足够真实”的吸引力，是下一章——第 22 章的主题。而水把这种吸引力玩到了极致，玩出了氢键、浮冰和地球的气候，那是第 23 章的舞台。至于手性这座迷宫，第 42 章会带你走到底：为什么你身体里的氨基酸几乎全是“左手”，而糖几乎全是“右手”。

\begin{tiaozhan}
\textbf{杂化轨道：一套成功的“记账方法”（跳过不影响主线）}

碳原子的电子排布是 $1s^2 2s^2 2p^2$，按理说只有 2 个 p 电子“落单”，凭什么形成 4 个等价的键？量子力学给出的描述是：成键时，1 个 2s 轨道和 3 个 2p 轨道的数学表达式可以重新组合成 4 个完全等价的“$\mathrm{sp^3}$ 杂化轨道”，它们自动指向正四面体的四个顶点——VSEPR 用排斥“猜”到的角度，杂化理论用波函数“算”了出来，两种模型互相印证。类似地，$\mathrm{sp^2}$ 杂化给出平面三角形（解释双键周围的 $120^{\circ}$），sp 杂化给出直线（解释三键）。

但要记住本章边界层的警告：实验测到的是电子密度的总分布，s、p、杂化都是我们切分这份分布的数学方式。把杂化轨道当成“分子内部真实存在的小房间”，就像把地图上的经纬线当成地球表面真实的绳子。模型因预言准确而被使用，不因被使用而变成实物。
\end{tiaozhan}

\section*{练一练}

\begin{sikaoti}
\item 氨 \chem{NH_3} 和水 \chem{H_2O} 的中心原子周围都是 4 对电子。画出两者的电子对排布示意图（成键对和孤对用不同标记），并解释为什么氨是三角锥形、水是折线形。图旁注明：这是表示法，电子对不是小珠子。
\item 三氟化硼 \chem{BF_3} 和氨 \chem{NH_3} 都含 4 个原子，为什么一个是平面三角形、一个是三角锥形？从中心原子的电子对总数和孤对数目回答。
\item 四氯化碳 \chem{CCl_4} 的 C–Cl 键是极性键，但 \chem{CCl_4} 整体不是极性分子；把其中一个氯换成氢，得到 \chem{CHCl_3}，就有了偶极。用四面体形状解释这一换一身的差别。
\item 二氧化硫 \chem{SO_2} 的中心硫原子周围有 3 团电子云（含 1 对孤对电子）。先预测它的形状，再判断它是不是极性分子，并说明你的依据。大气中的 \chem{SO_2} 能溶于水形成酸雨——你的判断与此相符吗？
\item 你的身体只能利用“左手”型的氨基酸来制造蛋白质。画一张信息流程图：从“分子的镜像形状”出发，经过“受体或酶的三维口袋”，到“能否被身体使用”。再用这张图解释：为什么橙子味和柠檬味的柠檬烯化学式完全相同。
\end{sikaoti}
